\section{Erd\H{o}s Problem 46: monochromatic unit fractions and rational targets}

\subsection{1) Problem and the theorem being used}
The question asks whether every finite colouring of the positive
integers has distinct monochromatic denominators
$2\leq n_1<\cdots<n_t$ with
\[
 \frac1{n_1}+\cdots+\frac1{n_t}=1.
\]
The problem page checked on 23 September 2026 records the affirmative
answer due to Croot. His finite theorem states that some absolute
$b>1$ has the following property: every $r$-colouring of
$\{2,\ldots,\lfloor b^r\rfloor\}$ contains a monochromatic finite
set $S$ whose reciprocal sum is $1$.

This note takes that precisely stated theorem as an external input.
It proves the passage to an arbitrary colouring, the infinitely many
disjoint representations, and the rational-target strengthening with
the colour and distinctness quantifiers checked. The substantial
analytic proof in Croot's paper is not repeated, and no new solution
or independent formal verification is claimed.

\subsection{2) The original assertion}
Let $\chi:\{2,3,\ldots\}\to\{1,\ldots,r\}$ be the given colouring.
Its restriction to Croot's finite interval supplies the required set
$S$. Ordering its members increasingly gives the denominators in the
question. As $S$ is a set, every denominator appears once. It contains
at least two elements because $1/n<1$ for every $n\geq2$.

\subsection{3) Moving every denominator beyond any fixed cutoff}
Fix an integer $T\geq2$. Give the integers $2,\ldots,T$ individual
fresh colours, all different from one another and from the original
$r$ colours. Retain $\chi$ on the integers larger than $T$.
The resulting colouring has $r+T-1$ colours, so Croot's theorem applies
again. A monochromatic reciprocal sum equal to $1$ cannot use a fresh
colour: its entire colour class is a singleton $\{n\}$ with $n\geq2$.
Therefore the resulting unit representation lies entirely beyond $T$
and is monochromatic for the original colouring.

Beginning with one such representation, choose the next cutoff larger
than all denominators used so far, and repeat. This produces infinitely
many pairwise disjoint finite sets $S_1,S_2,\ldots$, each of reciprocal
sum $1$, each monochromatic. Some fixed original colour occurs for
infinitely many of them, by the infinite pigeonhole principle.
Thus there is one colour containing infinitely many pairwise disjoint
unit representations. This last pigeonhole step is needed before
combining several representations into a larger rational target.

\subsection{4) Every positive rational has a monochromatic representation}
Let $a,b$ be positive integers. Colour $m\geq2$ by
\[
 \chi_b(m)=\chi(bm).
\]
The preceding argument applied to $\chi_b$ gives infinitely many
pairwise disjoint unit representations in a common induced colour.
Choose any $a$ of their denominator sets, say $S_1,\ldots,S_a$.
For each of them,
\[
 \sum_{m\in S_j}\frac1{bm}=\frac1b.
\]
The union
\[
 U=\bigcup_{j=1}^a bS_j
\]
therefore consists of distinct denominators, all of the same original
colour, and
\[
 \sum_{n\in U}\frac1n=\frac ab.
\]
Multiplication by $b$ is injective, and the sets $S_j$ are disjoint,
so denominators from different representations cannot collide.
Choosing the $S_j$ far enough along the construction makes every
denominator in $U$ exceed any prescribed cutoff.

In fact, for each fixed $b$ the same induced colour works for every
positive numerator $a$, because the infinite collection of unit
representations was chosen before specifying how many to combine.
We do not assert that the colour can be fixed independently of $b$.

\subsection{5) Checks and outcome}
Simply scaling a single unit representation by $b$ gives $1/b$, not
$a/b$. Taking $a$ copies with the same denominators would violate
distinctness, and choosing $a$ arbitrary monochromatic solutions
would not guarantee a common colour. The disjoint infinite family
and subsequent pigeonhole selection address both issues.

\textbf{Outcome: SOLVED, as a reconstruction using Croot's theorem.}
The original assertion and the stated disjointness and rational-target
consequences are fully deduced. The finite theorem remains the
explicitly cited external mathematical dependency.

\subsection{6) Sources}
\begin{itemize}
\item E. S. Croot III, \emph{On a coloring conjecture about unit
fractions}, Annals of Mathematics 157 (2003), 545--556, introductory
corollary, \texttt{https://annals.math.princeton.edu/wp-content/uploads/annals-v157-n2-p04.pdf}.
\item T. F. Bloom, Problem 46,
\texttt{https://www.erdosproblems.com/46}, checked 23 September 2026.
This page records both the disjoint representations and the rational
extension, which the above argument spells out.
\end{itemize}

The estimate measures coverage using the named published theorem, not
the originality of the result or an independent reproof of Croot's work.
\subsection{7) COMPLETION ESTIMATE}
\textbf{COMPLETION: 100\%}
